% ============================================================
\chapter{Étude et réalisation du Sprint 3}
\label{chap:sprint3}
% ============================================================

\section{Introduction}
\label{sec:intro_sprint3}

Dans ce chapitre, nous présentons l'étude et la réalisation du
\textbf{troisième sprint} de la plateforme FinBank. Après la mise en
place du backend et de la sécurité au Sprint 2, ce sprint est consacré
aux principales opérations bancaires proposées à l'utilisateur.

Le travail réalisé concerne la consultation du compte, les dépôts, les
retraits et les virements, ainsi que la consultation de l'historique et
le téléchargement des reçus. Il comprend également la gestion des
cartes, des investissements et des clés de paiement.

Le chapitre présente d'abord le backlog du sprint, puis la conception
des fonctionnalités à travers les diagrammes UML. Enfin, il décrit la
réalisation des services bancaires et prévoit des emplacements pour les
captures d'écran de l'application et de la documentation API.

\section{Backlog du Sprint}
\label{sec:backlog_sprint3}

Le tableau~\ref{tab:backlog_sprint3} présente les principales tâches
planifiées pour le Sprint 3.

\begin{table}[htbp]
\centering
\caption{Backlog du Sprint 3}
\label{tab:backlog_sprint3}
\renewcommand{\arraystretch}{1.3}
\begin{tabular}{|c|p{10cm}|c|}
\hline
\textbf{ID} & \textbf{Tâches} & \textbf{Estimation} \\
\hline
S3-T1 & Consulter le compte et afficher le solde disponible & 2 jours \\
\hline
S3-T2 & Implémenter les dépôts et les retraits & 3 jours \\
\hline
S3-T3 & Implémenter les virements entre comptes avec contrôle du solde & 4 jours \\
\hline
S3-T4 & Enregistrer les transactions et consulter l'historique & 2 jours \\
\hline
S3-T5 & Générer et télécharger les reçus des transactions & 2 jours \\
\hline
S3-T6 & Développer la gestion des cartes bancaires & 3 jours \\
\hline
S3-T7 & Développer les investissements et leur remboursement & 3 jours \\
\hline
S3-T8 & Gérer les clés de paiement et les virements par clé & 2 jours \\
\hline
S3-T9 & Vérifier les règles métier et documenter les endpoints & 2 jours \\
\hline
\multicolumn{2}{|r|}{\textbf{Total}} & \textbf{23 jours} \\
\hline
\end{tabular}
\end{table}

\section{Conception}
\label{sec:conception_sprint3}

La conception du Sprint 3 décrit les interactions de l'utilisateur
authentifié avec les services bancaires. Les opérations sont réalisées
par le backend en respectant les règles de propriété des comptes, de
solde disponible et de validation des données.

\subsection{Diagramme de cas d'utilisation}
\label{subsec:usecase_sprint3}

Le diagramme de cas d'utilisation regroupe les fonctionnalités
bancaires accessibles à l'utilisateur : consulter son compte, effectuer
des opérations, gérer ses moyens de paiement, investir et consulter
ses opérations passées.

\begin{figure}[htbp]
\centering
\IfFileExists{Figures/usecase_sprint3.png}{%
  \includegraphics[width=\textwidth,height=0.78\textheight,keepaspectratio]{Figures/usecase_sprint3.png}%
}{%
  \fbox{\parbox[c][5cm][c]{0.88\textwidth}{\centering
  Espace réservé au diagramme de cas d'utilisation du Sprint 3.}}
}
\caption{Diagramme de cas d'utilisation du Sprint 3}
\label{fig:usecase_sprint3}
\end{figure}

\subsection{Diagramme de classes du Sprint 3}
\label{subsec:classes_sprint3}

Le diagramme de classes présente les entités métier principales de ce
sprint. Un utilisateur possède un ou plusieurs comptes, transactions,
cartes, investissements et clés de paiement. Une transaction peut
associer un compte source et un compte cible, selon son type.

\begin{figure}[htbp]
\centering
\IfFileExists{Figures/classes_sprint3.png}{%
  \includegraphics[width=\textwidth,height=0.82\textheight,keepaspectratio]{Figures/classes_sprint3.png}%
}{%
  \fbox{\parbox[c][5cm][c]{0.88\textwidth}{\centering
  Espace réservé au diagramme de classes du Sprint 3.}}
}
\caption{Diagramme de classes du Sprint 3}
\label{fig:classes_sprint3}
\end{figure}

\subsection{Description textuelle : Consulter son compte}
\label{subsec:desc_compte_sprint3}

\begin{table}[htbp]
\centering
\caption{Cas d'utilisation : Consulter son compte}
\label{tab:uc_compte_sprint3}
\renewcommand{\arraystretch}{1.3}
\begin{tabular}{|p{3.2cm}|p{10.5cm}|}
\hline
\textbf{Élément} & \textbf{Description} \\
\hline
Titre & Consulter son compte \\
\hline
Acteur & Utilisateur (\texttt{USER}) \\
\hline
But & Consulter le numéro du compte, son solde et son état. \\
\hline
Précondition & L'utilisateur est authentifié et possède un compte. \\
\hline
Postcondition & Les informations actualisées du compte sont affichées. \\
\hline
Scénario principal &
1. L'utilisateur ouvre son espace bancaire. \newline
2. Le système identifie l'utilisateur authentifié. \newline
3. Le système recherche son compte. \newline
4. Le système retourne le solde et les informations du compte. \\
\hline
Scénario alternatif & Si le compte est introuvable ou inactif, le système retourne un message explicite. \\
\hline
\end{tabular}
\end{table}

\subsection{Description textuelle : Effectuer un virement}
\label{subsec:desc_virement_sprint3}

\begin{table}[htbp]
\centering
\caption{Cas d'utilisation : Effectuer un virement}
\label{tab:uc_virement_sprint3}
\renewcommand{\arraystretch}{1.3}
\begin{tabular}{|p{3.2cm}|p{10.5cm}|}
\hline
\textbf{Élément} & \textbf{Description} \\
\hline
Titre & Effectuer un virement \\
\hline
Acteur & Utilisateur (\texttt{USER}) \\
\hline
But & Transférer un montant vers un compte destinataire. \\
\hline
Précondition & L'utilisateur est authentifié, le montant est positif et les comptes sont valides. \\
\hline
Postcondition & La transaction est enregistrée avec son statut et les soldes sont mis à jour lorsque l'opération est validée. \\
\hline
Scénario principal &
1. L'utilisateur saisit le compte cible et le montant. \newline
2. Le système vérifie les données et le solde disponible. \newline
3. Le système débite le compte source et crédite le compte cible. \newline
4. Le système enregistre la transaction. \newline
5. Le système retourne le résultat de l'opération. \\
\hline
Scénario alternatif & En cas de solde insuffisant, de compte invalide ou de validation nécessaire, la transaction est refusée ou placée dans l'état approprié. \\
\hline
\end{tabular}
\end{table}

\subsection{Diagramme de classes participantes : Opération bancaire}
\label{subsec:participants_banking_sprint3}

Ce diagramme met en évidence l'interface, le contrôleur, le service,
les repositories et les entités mobilisés lors d'une opération bancaire.

\begin{figure}[htbp]
\centering
\IfFileExists{Figures/participants_operation_bancaire.png}{%
  \includegraphics[width=\textwidth,height=0.78\textheight,keepaspectratio]{Figures/participants_operation_bancaire.png}%
}{%
  \fbox{\parbox[c][5cm][c]{0.88\textwidth}{\centering
  Espace réservé au diagramme de classes participantes d'une opération bancaire.}}
}
\caption{Diagramme de classes participantes -- Opération bancaire}
\label{fig:participants_operation_sprint3}
\end{figure}

\subsection{Diagramme de séquence : Effectuer un virement}
\label{subsec:seq_virement_sprint3}

Le diagramme de séquence décrit le traitement d'un virement depuis la
requête de l'utilisateur jusqu'à l'enregistrement de la transaction.

\begin{figure}[htbp]
\centering
\IfFileExists{Figures/seq_virement.png}{%
  \includegraphics[width=\textwidth,height=0.82\textheight,keepaspectratio]{Figures/seq_virement.png}%
}{%
  \fbox{\parbox[c][5cm][c]{0.88\textwidth}{\centering
  Espace réservé au diagramme de séquence d'un virement.}}
}
\caption{Diagramme de séquence -- Effectuer un virement}
\label{fig:seq_virement_sprint3}
\end{figure}

\section{Réalisation}
\label{sec:realisation_sprint3}

Cette section présente les fonctionnalités implémentées durant le
Sprint 3 et leur intégration dans l'API REST de FinBank.

\subsection{Gestion des comptes et des transactions}
\label{subsec:comptes_transactions_sprint3}

Le contrôleur \texttt{AccountController} expose la consultation du
compte de l'utilisateur authentifié. Les opérations de dépôt, de
retrait et de virement sont regroupées dans le
\texttt{TransactionController}. Chaque requête est associée à
l'utilisateur issu du contexte de sécurité et non à un identifiant
fourni librement par le client.

Le service métier vérifie notamment la validité du montant, l'état des
comptes et la disponibilité du solde. Après le traitement, une entité
\texttt{Transaction} est créée avec son type et son statut. L'historique
permet ensuite de retrouver les opérations de l'utilisateur.

\begin{figure}[htbp]
\centering
\IfFileExists{Figures/swagger_transactions_sprint3.png}{%
  \includegraphics[width=0.9\textwidth]{Figures/swagger_transactions_sprint3.png}%
}{%
  \fbox{\parbox[c][4cm][c]{0.85\textwidth}{\centering
  Espace réservé à la capture Swagger des comptes et transactions.}}
}
\caption{Documentation Swagger des comptes et transactions}
\label{fig:swagger_transactions_sprint3}
\end{figure}

\subsection{Historique et téléchargement des reçus}
\label{subsec:historique_recus_sprint3}

L'utilisateur peut consulter la liste de ses transactions à partir de
l'endpoint d'historique. Il peut également télécharger le reçu PDF
d'une transaction. Avant de produire le document, le backend vérifie
que l'utilisateur est propriétaire du compte source ou du compte cible
associé à l'opération.

\begin{figure}[htbp]
\centering
\IfFileExists{Figures/history_receipt_sprint3.png}{%
  \includegraphics[width=0.9\textwidth]{Figures/history_receipt_sprint3.png}%
}{%
  \fbox{\parbox[c][4cm][c]{0.85\textwidth}{\centering
  Espace réservé à la capture de l'historique et du téléchargement d'un reçu.}}
}
\caption{Historique des transactions et téléchargement d'un reçu}
\label{fig:history_receipt_sprint3}
\end{figure}

\subsection{Gestion des cartes et des investissements}
\label{subsec:cartes_investissements_sprint3}

Le module des cartes permet à l'utilisateur de consulter ses cartes,
d'en créer une, de modifier ses informations autorisées et de la
supprimer. Le module des investissements permet de consulter les
investissements, d'investir un montant et de demander le remboursement
d'un investissement existant.

\begin{figure}[htbp]
\centering
\IfFileExists{Figures/cards_investments_sprint3.png}{%
  \includegraphics[width=0.9\textwidth]{Figures/cards_investments_sprint3.png}%
}{%
  \fbox{\parbox[c][4cm][c]{0.85\textwidth}{\centering
  Espace réservé aux captures des cartes et des investissements.}}
}
\caption{Interfaces de gestion des cartes et des investissements}
\label{fig:cards_investments_sprint3}
\end{figure}

\subsection{Clés de paiement}
\label{subsec:cles_paiement_sprint3}

Le module des clés de paiement permet d'enregistrer et de consulter les
clés associées à l'utilisateur. Une clé peut également être recherchée
afin d'identifier un destinataire avant l'exécution d'un transfert.
Les opérations restent protégées par l'authentification et les
validations métier du backend.

\begin{figure}[htbp]
\centering
\IfFileExists{Figures/payment_keys_sprint3.png}{%
  \includegraphics[width=0.9\textwidth]{Figures/payment_keys_sprint3.png}%
}{%
  \fbox{\parbox[c][4cm][c]{0.85\textwidth}{\centering
  Espace réservé à la capture de la gestion des clés de paiement.}}
}
\caption{Gestion des clés de paiement}
\label{fig:payment_keys_sprint3}
\end{figure}

\subsection{Vérification de la réalisation}
\label{subsec:verification_sprint3}

Les fonctionnalités du sprint sont accessibles à travers les
contrôleurs \texttt{/accounts}, \texttt{/transactions},
\texttt{/cards}, \texttt{/investments} et \texttt{/pix}. Les réponses
retournées par l'API utilisent des DTO afin de séparer les données
exposées des entités persistantes. Les erreurs de validation, de
propriété et de solde sont retournées à l'utilisateur sous forme de
réponses HTTP adaptées.

\section{Conclusion}
\label{sec:conclusion_sprint3}

Ce chapitre a présenté l'étude et la réalisation du Sprint 3 de
FinBank. La conception a été décrite à travers le diagramme de cas
d'utilisation, le diagramme de classes, le diagramme de classes
participantes et le diagramme de séquence du virement.

La réalisation de ce sprint a permis d'ajouter les fonctionnalités
bancaires opérationnelles : consultation des comptes, dépôts, retraits,
virements, historique, reçus, cartes, investissements et clés de
paiement. Ces services constituent le socle fonctionnel utilisé par les
interfaces et les processus présentés dans les chapitres suivants.
